% TOP_PROBLEM: {"id": 3415, "title": "Fundamental group torsion for subsets of Euclidean 3-space", "category_id": 7, "category": {"id": 7, "name": "topology", "display_name": "Topology", "description": "Properties preserved under continuous deformations.", "slug": "topology", "order_index": 7, "created_at": "2026-07-31T15:26:25.671Z"}, "status": "open", "problem_number": "OPG-37151", "set_id": 12}
% FIRST_POSTED: 2026-10-02
% ATTEMPT_STATUS: unresolved
% UnsolvedMath ID 3415; source catalogue code OPG-37151
% !TeX program = lualatex
% GPT-6 Astra Ultra research attempt, 2026-10-02; independently checked partial work.
% Completion estimate: no numerical percentage assigned; see final status and literature audit.
\documentclass[11pt,a4paper]{article}
\usepackage[margin=25mm]{geometry}
\usepackage{fontspec}
\setmainfont{lmroman10-regular.otf}[BoldFont=lmroman10-bold.otf,
 ItalicFont=lmroman10-italic.otf,BoldItalicFont=lmroman10-bolditalic.otf]
\usepackage{amsmath,amssymb,mathtools}
\usepackage{unicode-math}
\setmathfont{latinmodern-math.otf}
\AtBeginDocument{\let\setminus\smallsetminus}
\usepackage{xurl,hyperref}
\hypersetup{colorlinks=true,urlcolor=blue}
\setcounter{secnumdepth}{0}
\setlength{\parindent}{0pt}
\setlength{\parskip}{5pt}
\setlength{\emergencystretch}{3em}
\begin{document}
\section*{Fundamental group torsion for subsets of Euclidean 3-space}\label{3415}
{\small UnsolvedMath ID 3415 · OPG-37151 · Topology · OpenGarden}
\subsection{Definitions and mathematical statement}
For a subset $X\subseteq\mathbb R^3$ with the subspace topology and
a chosen point $x_0\in X$, the fundamental group $\pi_1(X,x_0)$
consists of continuous loops based at $x_0$, modulo homotopies that
fix the base point. Multiplication is concatenation. A nonidentity
element has finite order if some positive power of it is the identity.
The question is whether there exist $X,x_0$ and $[\gamma]\in\pi_1(X,x_0)$
such that
\[
                     [\gamma]\ne1,\qquad [\gamma]^m=1
                        \quad\text{for some integer }m\ge2.
\]
Equivalently, can a loop in a subset of three-space fail to contract
there while its $m$-fold concatenation does contract there?

No manifold, local contractibility, compactness, or regular-neighborhood
assumption is imposed on $X$. One may restrict to the path component
containing $x_0$, since based loops and their homotopies remain in
that component. Thus wild embedded subsets are within the scope.
The identity element, whose order is one, is excluded; otherwise
the literal phrase ``an element of finite order'' would make the
question trivial. A positive answer requires both a null-homotopy
for a nontrivial power and a proof that the original loop is not
null-homotopic in the same subset.
\subsection{Short English statement}
Can a possibly wild subset of three-dimensional Euclidean space have a nontrivial finite-order element in its fundamental group?
\subsection{Research attempt}
\paragraph{Scope and process.}
GPT-6 Astra Ultra, 2 October 2026. The canonical formulation above is retained. This attempt uses a primary-source literature audit, independent restricted proofs, and adversarial checks. Reproved results are not claimed as new. Numerical tests below are reproducible in \texttt{scripts/check\_attempt\_batch03\_extremal\_2026\_10\_02.py}.

\paragraph{Scope and literature audit.}
The original entry [S2] asks about arbitrary subsets, including wild compact sets, and attributes the problem to folklore. Tame neighbourhood or planar results cannot simply be applied to all such subsets. We avoid transferring a singular-homology conclusion from Alexander duality without the necessary hypotheses. Instead, the attempt proves a compact-witness reduction and a precise conditional neighbourhood criterion, as well as excluding the class of spaces with finite graph spines.

\paragraph{Every positive answer has a compact witness.}
Suppose $X\subseteq\mathbb R^3$ has a based loop $\gamma$ representing a nonidentity element $g$ with $g^m=1$ for some integer $m\ge2$. Choose a nullhomotopy of the repeated loop $\gamma^m$, expressed as a continuous map $F:D^2\to X$ with the prescribed boundary loop. Set $K=F(D^2)$. Since $D^2$ is compact and path connected, $K$ is a compact path-connected subset of $\mathbb R^3$. Its boundary image contains the whole image of $\gamma$, so $\gamma$ is also a loop in $K$.

The same map $F$ proves $[\gamma]^m=1$ in $\pi_1(K)$. Yet $[\gamma]\ne1$ there: a nullhomotopy in $K$ would also be a nullhomotopy in $X$ by composing with inclusion, contradicting the choice of $g$. Thus a positive answer for arbitrary subsets gives a positive answer for compact path-connected subsets. The converse is immediate by inclusion among the permitted subsets. Notice that this proves compactness of a witness, not local contractibility or existence of a regular neighbourhood. A continuous image of a disk need not be a disk or a manifold.

\paragraph{Excluding finite graph spines.}
Suppose a path component of $X$ deformation retracts to a finite connected graph $Y$. Retraction induces an isomorphism of fundamental groups. Collapsing a maximal tree of $Y$ gives a wedge of circles, so $\pi_1(Y)$ is the free group on its non-tree edges. To check torsion-freeness directly, let a nonidentity element have reduced word $w$. Removing inverse first and last letters repeatedly writes it as $w=sus^{-1}$ with $u$ a nonempty cyclically reduced word. The first and last letters of $u$ are not inverse. Hence $u^m$ is still a nonempty reduced word for every positive $m$. The reduced form of $w^m$ is $su^ms^{-1}$, which is nonempty. No nonidentity element has finite order.

This proves a negative answer for every subset having such a deformation retraction, regardless of how its graph spine is embedded. It does not assert that every compact subset of three-space has a finite graph spine. In particular the compact-image reduction above cannot be followed by an unjustified finite-complex replacement.

\paragraph{The exact neighbourhood-detection gap.}
Suppose a compact path-connected $K$ with basepoint $x_0$ lies in nested path-connected neighbourhoods $U_1\supseteq U_2\supseteq\cdots$, each having torsion-free fundamental group. Inclusion induces a homomorphism
\[
 \Phi:\pi_1(K,x_0)\longrightarrow
 \varprojlim_j\pi_1(U_j,x_0).
\]
If this homomorphism is injective, then $\pi_1(K,x_0)$ is torsion-free. Indeed, if $g^m=1$, every coordinate of $\Phi(g)$ has finite order dividing $m$ and must be the identity. Thus $\Phi(g)=1$, and injectivity forces $g=1$.

The logical burden is injectivity. The assertion that a loop contracts in every neighbourhood of $K$ only supplies a sequence of possibly unrelated disks. It does not supply a limiting continuous disk whose image lies in $K$. Compactness of the target alone gives no equicontinuity or compatibility of these maps. Our criterion is therefore conditional and cannot be applied to arbitrary compact sets by simply taking smaller neighbourhoods.

\paragraph{Verification and final status.}
All maps in the compact reduction preserve the chosen basepoint, and nontriviality is checked in the correct direction under inclusion. The script checks free-word powers for small reduced words as a diagnostic of the graph-spine argument; the proof covers every length. The first unproved general step is a torsion-detection theorem for compact wild subsets, or an explicit compact subset with a finite-order essential loop. Neither is supplied. The arbitrary-subset existence question is \textbf{UNRESOLVED} by this attempt.

\subsection{Sources}
\begin{itemize}
\item[S1] ULAM.AI and the UnsolvedMath Contributors, supplied problems.json, record 3415 (OPG-37151), OpenGarden collection. Catalogue identity and source transcription; CC BY 4.0.
\url{https://huggingface.co/datasets/ulamai/UnsolvedMath}.
\item[S2] Open Problem Garden, Fundamental group torsion for subsets of Euclidean 3-space. Original problem and discussion; the mathematical formulation below is expanded from the supplied source record.
\url{http://www.openproblemgarden.org/op/torsion_for_subsets_of_mathbb_r_3}.

\end{itemize}
\end{document}
